\section{Evidencia numérica}\label{sec:numerico}

Nobel, Agrawal y Boyd~\cite{NAB} dan $\alpha\in[1{,}944000752,\,1{,}944001082]$.
Para las dos tablas usamos como referencia $\alpha_0=1{,}944001$ y $\Phi_0=1-\alpha_0$,
porque $\alpha=-s_1$ (Sección~\ref{sec:preliminares}). Las extrapolaciones carecen de cota de error
certificada, y la forma de los ajustes, en potencias de $1/n$, es una hipótesis de trabajo: el
Teorema~\ref{thm:intro-constante} da el límite, pero no la velocidad de convergencia.

\paragraph{El predictor.}
Pongamos $C'_{{\rm pred},0}(n)=\log Q_{\rm pred}(n)-\log n!-n\Phi_0$. Una resolución numérica del
maximizador
canónico encuentra interioridad en los diez tamaños tabulados, en aritmética de punto flotante.
El pseudodeterminante de $S_n$ completa el cálculo de la Tabla~\ref{tab:predictor}.
Con $\alpha_0$ y $40\le n\le240$, el ajuste $C+a/n+b/n^2$ da $C=0{,}94015$ y el que añade $c/n^3$ da
$C=0{,}94013$.
En el intervalo completo de $\alpha$, los ajustes que empiezan en $40$ o $60$ dan extremos redondeados de
$0{,}9400$ y $0{,}9402$.
Al incluir $20$ y $27$ los extremos redondeados son $0{,}9397$ y $0{,}9404$. Así se obtiene la estimación
$C_{\rm pred}\approx0{,}9401$.

\begin{table}[ht]
\centering
\begin{tabular}{lccccc}
\toprule
$n$ & 20 & 27 & 40 & 60 & 80\\
$C'_{{\rm pred},0}(n)$ & 0,9112 & 0,9165 & 0,9239 & 0,9290 & 0,9316\\
\midrule
$n$ & 100 & 130 & 160 & 200 & 240\\
$C'_{{\rm pred},0}(n)$ & 0,9333 & 0,9348 & 0,9358 & 0,9367 & 0,9372\\
\bottomrule
\end{tabular}
\caption{Término constante del predictor para tableros finitos.}\label{tab:predictor}
\end{table}

\paragraph{Evaluación directa de la fórmula.}
Por separado, calculamos el perfil continuo $\nu_1$ resolviendo numéricamente sus ecuaciones de masa uno, y
con él $I_2$ y $\det_2(I+J)$ mediante una discretización de Nyström~\cite{Bornemann} del operador de la
Sección~\ref{sec:D03}, en mallas sucesivas y con extrapolación en el paso. Se obtiene
$I_2\approx1{,}02638534$ y $\det_2(I+J)\approx0{,}005162694$, y la fórmula del
Teorema~\ref{thm:intro-constante} da $C_{\rm pred}\approx0{,}940144$, dentro de las bandas anteriores. Este
cálculo no usa $\alpha$ ni la Tabla~\ref{tab:predictor}, y tampoco tiene cota de error certificada.
En las parejas de mallas de tamaños $192$ y $384$ (pares), y $193$ y $385$ (impares), las extrapolaciones
dieron, respectivamente, $0{,}94014395698$ y $0{,}94014395711$.

\paragraph{El conteo exacto.}
Los valores exactos de $Q(n)$ se conocen hasta $n=27$~\cite{OEIS,PE} desde 2016; el de $n=27$ se
confirmó en 2025 con otro programa~\cite{YaoLi}. Con los valores de
$16\le n\le27$, $C'_0(n)=\log Q(n)-\log n!-n\Phi_0$ toma los valores de la Tabla~\ref{tab:exacto}.

\begin{table}[ht]
\centering
\begin{tabular}{lcccccc}
\toprule
$n$ & 16 & 17 & 18 & 19 & 20 & 21\\
$C'_0(n)$ & 0,9404 & 0,9209 & 0,9135 & 0,9224 & 0,9320 & 0,9187\\
\midrule
$n$ & 22 & 23 & 24 & 25 & 26 & 27\\
$C'_0(n)$ & 0,9178 & 0,9241 & 0,9295 & 0,9272 & 0,9265 & 0,9285\\
\bottomrule
\end{tabular}
\caption{Término constante del conteo exacto.}\label{tab:exacto}
\end{table}

En noviembre de 2026 se publicó, en una prepublicación, el valor de $Q(28)$~\cite{Cherepanov},
obtenido con un algoritmo nuevo de conteo exacto. Poco después la OEIS~\cite{OEIS} incorporó también
$Q(29)$, calculado por P.~A.~Webb con una implementación de la descomposición por colores que, según
esa entrada, se usó para $Q(28)$ y $Q(29)$. Hasta donde sabemos, ninguno de los dos valores se ha
comprobado con un algoritmo distinto, por lo que no los incluimos en la Tabla~\ref{tab:exacto} ni en
la Figura~\ref{fig:comparacion}. A título informativo, con el primero resulta
$C'_0(28)\approx0{,}9305$, dentro del intervalo de los valores de la tabla.

La Figura~\ref{fig:comparacion} permite leer las dos tablas en la misma escala.
Al aumentar $n$ los puntos se acercan al eje vertical porque la abscisa es $1/n$.

\begin{figure}[htbp]
\centering
\begin{tikzpicture}[x=158cm,y=103cm]
 \begin{scope}[shift={(0,-.905)}]
 \foreach \v/\letra in {.910/{0,91},.920/{0,92},.930/{0,93},.940/{0,94}}{
   \draw[black!18] (0,\v)--(.065,\v);
   \node[anchor=east,font=\footnotesize] at (-.0007,\v) {\letra};}
 \draw[btinta] (0,.905)--(.066,.905) (0,.905)--(0,.945);
 \foreach \v/\letra in {0/{0},.02/{0,02},.04/{0,04},.06/{0,06}}{
   \draw[btinta] (\v,.905)--(\v,.9045);
   \node[anchor=north,font=\footnotesize] at (\v,.904) {\letra};}
 \draw[btinta,dashed] (0,.940144)--(.065,.940144);
 \draw[btinta,line width=.75pt] plot coordinates {(0.050000000,0.911188472) (0.037037037,0.916522978) (0.025000000,0.923873790) (0.016666667,0.928966033) (0.012500000,0.931643280) (0.010000000,0.933288384) (0.007692308,0.934831567) (0.006250000,0.935807595) (0.005000000,0.936660345) (0.004166667,0.937232189)};
 \foreach \x/\y in {0.050000000/0.911188472,0.037037037/0.916522978,0.025000000/0.923873790,0.016666667/0.928966033,0.012500000/0.931643280,0.010000000/0.933288384,0.007692308/0.934831567,0.006250000/0.935807595,0.005000000/0.936660345,0.004166667/0.937232189}{\fill[btinta] (\x,\y) circle[radius=1.45pt];}
 \foreach \x/\y in {0.062500000/0.940434608,0.058823529/0.920874442,0.055555556/0.913509083,0.052631579/0.922429638,0.050000000/0.931979175,0.047619048/0.918660405,0.045454545/0.917777869,0.043478261/0.924122583,0.041666667/0.929528249,0.040000000/0.927235036,0.038461538/0.926480702,0.037037037/0.928458576}{\draw[btinta,fill=white,line width=.65pt]
    (\x,\y) ++(-1.55pt,-1.55pt) rectangle ++(3.1pt,3.1pt);}
 \node[anchor=north,font=\small] at (.033,.9012) {$1/n$};
 \end{scope}
\end{tikzpicture}
\caption{Términos centrados con la misma referencia $\Phi_0$. Los círculos, unidos por segmentos,
representan el predictor de la Tabla~\ref{tab:predictor}; los cuadrados, sin unir, el conteo exacto
de la Tabla~\ref{tab:exacto}. La línea horizontal discontinua marca la evaluación de la constante
del predictor mediante la discretización de Nyström. Los tamaños representados son los de las dos tablas.}
\label{fig:comparacion}
\end{figure}


En $n=20$ y $n=27$ se tiene $\log Q(n)-\log Q_{\rm pred}(n)=0{,}0208$ y $0{,}0119$. Los dos números
son compatibles con una diferencia de orden $1/n$, pero no permiten estimar su límite.

Los programas y los datos con los que se obtuvieron las cifras de esta sección forman el paquete de
verificación numérica que acompaña al artículo~\cite{Codigo}.
